\section{Estimates for the terminal tail}\label{app:partial-turn-estimates}

\subsection{The initial-angle exclusion}

\begin{proposition}\label{app:partial-initial}
For every downward cap $U$ and $\pi/4\le\alpha\le\arctan(8/3)$,
\[
 P_\alpha(U)<\frac{5259}{6400}<\frac{411}{500}<\frac\Mc2.
\]
\end{proposition}
\begin{proof}
First put $\kappa=\cot\alpha$ and center the projection as
$I=[-2C,2C]$. Write $A=f_U$ and
$H=h_U(\alpha)/\sin\alpha$. The roof satisfies $A(x)\le H-\kappa x$
and the whole terminal wall is $H-\csc\alpha-\kappa x$.
The left exterior roof is at most one; integrate the support bound
on $[C,2C]$ and the terminal lower bound on $[-C,0]$, even where
that affine bound is negative. This gives
\[
 P_\alpha\le(1+\sqrt{1+\kappa^2})C-2\kappa C^2
 \le\frac{(1+\sqrt{1+\kappa^2})^2}{8\kappa}.
\]
On $4/5\le\kappa\le1$ the last expression decreases: its derivative
has the sign of $(\sqrt{1+\kappa^2}-2)(\sqrt{1+\kappa^2}+1)$.
Its largest value is $(33+5\sqrt{41})/80<1301/1600<41/50$.
This handles $\alpha\le\arctan(5/4)$.

\paragraph{Reduction and the three actual supports}

\Cref{lem:selection-canonical,lem:selection-maximum} supplies height extrusion, middle-chord reduction, coercivity, attainment, and top insertion. The attainment proof also applies with
\(\alpha\) fixed: if a cap has score greater than \(41/50\), the
fixed-angle maximum is attained inside the same compact width interval.
Indeed the width-six calculation in \cref{lem:selection-width}
is below $3/4$, so every cap with score above $41/50$ has
width below six. A maximizing sequence at this fixed angle can
therefore be centered in the same compact box; score continuity
and Blaschke compactness give its maximum.
Its canonical representative has height one, an affine middle roof, and
a top face meeting the middle window. Thus, a hypothetical cap with score
at least \(5259/6400\) yields such a canonical fixed-angle maximizer
with at least the same score. It suffices to rule out that maximizer.


Center the projection and introduce the parameters

\begin{equation}\label{eq:partial-at-2}
\begin{gathered}
I=[-2C,2C],\quad J=[-C,C],\\
K=\sqrt2-1,\quad \kappa=\cot\alpha,\quad
d=\frac{\kappa}{1+\sqrt{1+\kappa^2}}.
\end{gathered}
\end{equation}

The outgoing-wall estimate above handles
\(\alpha\le\arctan(5/4)\). For the new interval it suffices to take

\begin{equation}\label{eq:partial-at-3}
\begin{gathered}
\frac38\le\kappa\le\frac45,\qquad
\kappa d=\sqrt{1+\kappa^2}-1,\\
\kappa=\frac{2d}{1-d^2},\qquad 0<d<K.
\end{gathered}
\end{equation}

Use the actual supports to define

\begin{equation}\label{eq:partial-at-4}
u=\max_U(x+y)-1,\quad v=\max_U(-x+y)-1,\quad
b=\max_U\left(x+\frac{y-1}{\kappa}\right).
\end{equation}

The three outer support lines are
\(1+u-x\), \(1+v+x\), and \(1-\kappa(x-b)\). The actual outgoing
inner wall is \(\kappa(b-d-x)\). Consequently

\begin{equation}\label{eq:partial-at-5}
\begin{aligned}
A(x)&\le \min\{1,1+u-x,1+v+x,1-\kappa(x-b)\},\\
N_\alpha(U)(x)&\ge
\max\{0,\min(u-K-x,v-K+x),\kappa(b-d-x)\}.
\end{aligned}
\end{equation}

The 45-degree angle is genuinely visited throughout the new interval.
Only the displayed lower bound on the true niche is used.

There is an essential compatibility condition between the two first
supports:

\begin{equation}\label{eq:partial-at-6}
\boxed{b\le u\le 2(1-\kappa)C+\kappa b,
\qquad u,v,b\le2C.}
\end{equation}

Indeed, \(y\le1\) and \(\kappa<1\) give
\(x+(y-1)/\kappa\le x+y-1\), hence \(b\le u\). At a point attaining
\(u\),

\[
\kappa b\ge\kappa x+y-1=u-(1-\kappa)x
\ge u-2(1-\kappa)C.
\]

The remaining upper bounds follow directly from the containing box.

Since the middle is affine and its top meets \(J\), either
\(A(-C)=1\) or \(A(C)=1\). The first alternative gives
\(v\ge C\) and \(u,b\ge-C\); the second gives \(u,b\ge C\) and
\(v\ge-C\). These alternatives will be handled directly.

\paragraph{Two preliminary width cuts}

\paragraph{The 45-degree relaxation alone handles \(C\le49/100\)}

Put

\[
g_C(z)=\int_C^{2C}(x-z)_+\,dx
=\frac{(2C-z)_+^2-(C-z)_+^2}{2}.
\]

For this paragraph only, the 45-degree relaxation is reflection
invariant, so label the high endpoint as the left one. Then
\(v\in[C,2C]\), \(u\in[-C,2C]\), and the support bounds give

\begin{equation}\label{eq:partial-at-7}
P_{45}\le F_{45}(u,v):=2C-g_C(u)-g_C(v)
-\int_{-C}^C[\min(u-K-x,v-K+x)]_+\,dx.
\end{equation}

This function is concave on that rectangle. Here is a clipping-inclusive
verification which will also be used below. For an integrated maximum of
two affine lines with different spatial slopes, the moving intersection
contributes the nonnegative second variation

\begin{equation}\label{eq:partial-at-8}
\frac{(\hbox{difference of varied intercepts})^2}
{|\hbox{difference of spatial slopes}|}.
\end{equation}

An integrated minimum has the opposite sign. These formulas follow by
integrating the triangle swept out by the moving intersection. A switch
outside the integration interval contributes zero. The resulting
integrals are continuously differentiable and piecewise quadratic in
the intercepts, so checking their quadratic second variations on the
open pieces also checks concavity across all clipping boundaries.

For the loss \(2C-F_{45}\), the sole negative term can occur at an
exposed tent apex inside \(J\); it is
\(- (\xi-\eta)^2/2\), for variations \((\xi,\eta)\) of \((u,v)\).
The left outer switch supplies \(\eta^2\). If \(u>C\), the right
outer switch supplies \(\xi^2\). If \(u\le C\) and the apex is
positively exposed, its falling side meets the floor at
\(u-K\in(-C,C)\), supplying \(\xi^2\) instead. Thus the loss has
nonnegative second variation; cases without an exposed interior apex
already have only nonnegative terms. Boundary parameters follow by
continuity.

For \(C\ge K/2\), the common critical point
\(u=v=C+K/2\) lies in the rectangle. Its positive tent is contained in
\(J\), and \eqref{eq:partial-at-7} is at most

\begin{equation}\label{eq:partial-at-9}
2C-2(C-K/2)^2.
\end{equation}

For \(C<K/2\), the trivial bound \(P_{45}\le2C<K\) suffices. The
right side of \eqref{eq:partial-at-9} increases up to \(C=49/100\), and

\begin{equation}\label{eq:partial-at-10}
\frac{49}{50}-\frac12\left(\frac{49}{50}-K\right)^2
<\frac{41}{50}.
\end{equation}

For an exact check, \(\sqrt2<99/70\) gives
\(49/50-K=99/50-\sqrt2>2\sqrt2/5\). Hence every remaining
canonical cap has \(C>49/100\).

\paragraph{For \(\kappa\ge3/5\), the outgoing wall handles \(C\le1/2\)}

The outgoing-wall calculation at the start of this proof gives, for every cap,

\begin{equation}\label{eq:partial-at-11}
P_\alpha\le(1+\sqrt{1+\kappa^2})C-2\kappa C^2.
\end{equation}

For \(C\le1/2\) and \(3/5\le\kappa\le4/5\), this expression
increases with \(C\); at \(C=1/2\) it decreases with \(\kappa\).
Therefore it is at most

\begin{equation}\label{eq:partial-at-12}
\frac{2+\sqrt{34}}{10}<\frac45.
\end{equation}

It remains to consider

\begin{equation}\label{eq:partial-at-13}
\begin{cases}
C>49/100,&3/8\le\kappa\le3/5,\\
C>1/2,&3/5\le\kappa\le4/5.
\end{cases}
\end{equation}

\paragraph{The exact finite-support relaxation and the high-right case}

For \(b\le u\), write

\begin{equation}\label{eq:partial-at-14}
\begin{gathered}
x_0=\frac{u-\kappa b}{1-\kappa},\\
\begin{aligned}
Q_C(u,b)&=\int_C^{2C}\max\{0,x-u,\kappa(x-b)\}\,dx\\
&=\kappa g_C(b)+(1-\kappa)g_C(x_0).
\end{aligned}
\end{gathered}
\end{equation}

The identity follows from \(b\le u\le x_0\): the three successive
active pieces are zero, \(\kappa(x-b)\), and \(x-u\), with any
pieces outside \([C,2C]\) automatically clipped by \(g_C\).

Define

\begin{equation}\label{eq:partial-at-15}
\begin{aligned}
N(u,v,b)&=\int_{-C}^C
\max\{0,\min(u-K-x,v-K+x),\kappa(b-d-x)\}\,dx,\\
F(u,v,b)&=2C-g_C(v)-Q_C(u,b)-N(u,v,b).
\end{aligned}
\end{equation}

In either high-endpoint alternative the left exterior majorant in \eqref{eq:partial-at-5}
is \(\min(1,1+v+x)\), and the right exterior majorant is
\(1-\max(0,x-u,\kappa(x-b))\). Thus
\(P_\alpha\le F\).

Suppose first that the high endpoint is on the right and \(v<C\).
Increase \(v\) to \(C\), keeping \(u,b\) fixed. This is a variation
of the relaxation; it need not preserve realizability by a cap. The
charged left exterior reward increases at rate \(C\). The derivative
of \(N\) is the length of the interval on which the rising tent line
\(v-K+x\) is exposed. Put

\begin{equation}\label{eq:partial-at-16}
p=\frac{u-v}{2},\quad e=b-d,\quad
x_L=\frac{\kappa e-v+K}{1+\kappa}.
\end{equation}

Every such exposed point lies between \(x_L\) and \(p\), so its
length is at most \((p-x_L)_+\). Since \(u\le2C\), \(v\le C\),
and \(b\ge C\),

\begin{equation}\label{eq:partial-at-17}
p-x_L\le
\frac{(3-\kappa)C+2\kappa d-2K}{2(1+\kappa)}<C.
\end{equation}

The strict inequality uses \(\kappa>1/3\) and \(\kappa d<K\).
Hence \(F\) strictly increases until \(v=C\). All of \eqref{eq:partial-at-6} is
preserved. It is therefore enough to bound \(F\) on the convex domain

\begin{equation}\label{eq:partial-at-18}
\boxed{D_C=\{C\le v\le2C,\ -C\le b\le2C,
\ b\le u\le2(1-\kappa)C+\kappa b\}.}
\end{equation}

This also contains the original high-left alternative. No reflection of
the partial objective has been used.

\paragraph{Two convex domains cover the useful part of the relaxation}

Put

\begin{equation}\label{eq:partial-at-19}
a=u-K,\quad c=v-K,\quad e=b-d,\quad
Y_0=(1+\kappa)C+K-v+\kappa(b-d).
\end{equation}

Thus \(Y_0\ge0\) means \(x_L\ge-C\). Define the two convex sets

\begin{equation}\label{eq:partial-at-20}
D_1=D_C\cap\{a\le C\},\qquad
D_2=D_C\cap\{Y_0\ge0\}.
\end{equation}

\emph{Concavity claim.} The function $F$ is concave on each of $D_1,D_2$.

 Work at generic parameters in the relative interiors of the
finite quadratic pieces, and use variations \((\xi,\eta,\zeta)\)
of \((u,v,b)\). As in \eqref{eq:partial-at-8}, all switches of the positive niche
contribute nonnegative second variations except an exposed tent apex
\(p=(u-v)/2\) inside \(J\). If there is no such apex, the loss
\(2C-F=g_C(v)+Q_C+N\) is already convex on that piece. Otherwise the
only negative term is

\begin{equation}\label{eq:partial-at-21}
-\frac{(\xi-\eta)^2}{2}.
\end{equation}

Throughout the interior of \(D_C\), the left outer switch contributes
\(\eta^2\). The complete compensation is as follows.

\emph{On \(D_1\).} If \(a\ge e\), the exposed falling tent meets the
floor at \(x=a\in(-C,C)\). This contributes \(\xi^2\).
If \(a<e\), it instead meets the terminal wall at

\[
x_R=\frac{a-\kappa e}{1-\kappa}\in(p,C),
\]

which contributes \((\xi-\kappa\zeta)^2/(1-\kappa)\).
For \(e<C\), the subsequently exposed terminal wall meets the floor
at \(e\in(-C,C)\), contributing \(\kappa\zeta^2\). For
\(e>C\), instead \(b=e+d\in(C,2C)\), and the outer right
height/terminal switch at \(b\) supplies the same contribution.
In both cases

\begin{equation}\label{eq:partial-at-22}
\frac{(\xi-\kappa\zeta)^2}{1-\kappa}
+\kappa\zeta^2\ge\xi^2.
\end{equation}

Thus the positive contributions include \(\xi^2+\eta^2\), which
dominate \eqref{eq:partial-at-21}.

\emph{On \(D_2\), outside \(D_1\).} Now \(u>C+K\). Consequently
\(C<x_0<2C\), so the outer terminal/45-degree switch contributes
\((\xi-\kappa\zeta)^2/(1-\kappa)\).
If \(c+e>0\), the rising tent meets the terminal wall at
\(x_L\in(-C,p)\); this contributes
\((\eta-\kappa\zeta)^2/(1+\kappa)\). The inequality

\begin{equation}\label{eq:partial-at-23}
\frac{(\xi-\kappa\zeta)^2}{1-\kappa}
+\frac{(\eta-\kappa\zeta)^2}{1+\kappa}
\ge\frac{(\xi-\eta)^2}{2}
\end{equation}

cancels \eqref{eq:partial-at-21}.

If \(c+e<0\), the rising tent instead emerges from the floor at
\(-c>e\). The support compatibility and \(u>C+K\) give

\begin{equation}\label{eq:partial-at-24}
\kappa(e+C)\ge u-(2-3\kappa)C-\kappa d
>(3\kappa-1)C+K-\kappa d>0.
\end{equation}

Hence \(-C<e<-c<p\). The terminal/floor switch at \(e\) supplies
\(\kappa\zeta^2\), and the rising-tent/floor switch at \(-c\)
supplies \(\eta^2\). \eqref{eq:partial-at-22} again supplies \(\xi^2\), cancelling
the only negative term.

All remaining switch terms are nonnegative. The affine spatial slopes
\(0,1,-1,-\kappa\) are distinct, so ties have zero spatial measure;
the integrals in \eqref{eq:partial-at-15} are continuously differentiable and piecewise
quadratic in the parameters. The verified nonnegative second variations
therefore prove convexity of the loss on each convex set, including all
degenerate switch and window boundary cases by continuity. This proves the concavity claim.

\paragraph{A common critical point, with all its clipping checked}

Define

\begin{equation}\label{eq:partial-at-25}
\begin{aligned}
b_*&=\frac{K+d}{2},\\
u_*&=(1-\kappa)C+\frac K2+\frac{\kappa d}{2},\\
v_*&=(1+\kappa)C+\frac K2-\frac{\kappa d}{2}.
\end{aligned}
\end{equation}

Its relevant switch positions are

\begin{equation}\label{eq:partial-at-26}
\begin{aligned}
x_{0,*}&=C+K/2,&x_{L,*}&=K/2-C,\\
p_*&=-\kappa(C-d/2),&e_*&=(K-d)/2,\\
a_*&=u_*-K.&&
\end{aligned}
\end{equation}

The only potentially delicate clipping condition is \(a_*\ge e_*\).
The exact identities in \eqref{eq:partial-at-3} and \(K^2=1-2K\) give

\begin{equation}\label{eq:partial-at-27}
a_*-e_*=(1-\kappa)(C-1/2)
+\frac{(K-d)^2}{2(1+d)}.
\end{equation}

This is positive for \(C\ge1/2\). For
\(C>49/100\), \(3/8\le\kappa\le3/5\), use
\(K>41/100\), \(d<7/25\), and \(1-\kappa\le5/8\). They give

\begin{equation}\label{eq:partial-at-28}
a_*-e_*>-\frac1{160}+\frac{169}{25600}
=\frac9{25600}>0.
\end{equation}

Here \(d\le(\sqrt{34}-5)/3<7/25\), because
\(\sqrt{34}<146/25\). Thus \eqref{eq:partial-at-13} proves the required clipping in
every remaining case.

We have \(b_*<C\), \(C<x_{0,*}<2C\), and

\begin{equation}\label{eq:partial-at-29}
-C<x_{L,*}<p_*<0<e_*<a_*<C.
\end{equation}

For example \(p_*-x_{L,*}=a_*>0\),
\(a_*-p_*=C-K/2>0\), and \(C>d/2\); these identities check the
displayed ordering. Also \(2C-v_*=a_*>0\),
\(v_*-C=\kappa(C-d/2)+K/2>0\), and
\(u_*-b_*=(1-\kappa)(C-d/2)>0\). Together with the value of
\(x_{0,*}\), these show that the point is in the interior of
\(D_C\). It lies in both \(D_1,D_2\), since
\(a_*<C\) and \(Y_{0,*}=(1+\kappa)K/2>0\).

The outer first-45-degree, second-45-degree, and terminal horizontal
exposure lengths are respectively

\[
C-K/2,\qquad a_*,\qquad K/2.
\]

By \eqref{eq:partial-at-29} the corresponding niche exposure lengths are exactly the same:
\(a_*-p_*\), \(p_*-x_{L,*}\), and \(x_{L,*}+C\).
Consequently all three first derivatives of \(F\) vanish at \eqref{eq:partial-at-25}.
The concavity claim proves that this point maximizes \(F\) on each of
\(D_1,D_2\).

For the value, put \(z=C-K/2\), \(s=\kappa(C-d/2)\). Direct
integration over the displayed intervals gives

\[
\begin{aligned}
g_C(v_*)&=(z-s)^2/2,\\
Q_C(u_*,b_*)&=\kappa(3C^2/2-Cb_*)+(1-\kappa)z^2/2,\\
N(u_*,v_*,b_*)&=z^2-s^2/2+Ks/2+\kappa K^2/8.
\end{aligned}
\]

Their sum is \(2z^2+Ks\). Thus on \(D_1\cup D_2\),

\begin{equation}\label{eq:partial-at-30}
\boxed{F\le F_*:=2C-2(C-K/2)^2-\kappa K(C-d/2).}
\end{equation}

\paragraph{The complement is strictly below \(319/400\)}

Suppose \((u,v,b)\in D_C\setminus(D_1\cup D_2)\). Then
\(u>C+K\) and \(Y_0<0\). Combining the latter with \eqref{eq:partial-at-6} gives

\begin{equation}\label{eq:partial-at-31}
v-u>(3\kappa-1)C+K-\kappa d>0.
\end{equation}

In particular \(u,v>C+K\), so both ends of the 45-degree tent are
strictly positive. Its apex belongs to \((-C,C)\), since
\(u,v\in(C+K,2C]\). Dropping the outgoing support and charge only
enlarges the relaxation. With \(h=(u+v)/2\), exact 45-degree
integration therefore gives

\begin{equation}\label{eq:partial-at-32}
F\le 2C-2C^2+2KC-(h-C)^2.
\end{equation}

Moreover, \eqref{eq:partial-at-31} implies

\begin{equation}\label{eq:partial-at-33}
h-C>\frac{3K}{2}-\frac{\kappa d}{2}>\frac9{20},
\end{equation}

using \(K>2/5\) and
\(\kappa d\le\sqrt{41}/5-1<3/10\). Finally
\(1+K=\sqrt2\), so

\[
2(1+K)C-2C^2\le1.
\]

Equations \eqref{eq:partial-at-32}–\eqref{eq:partial-at-33} prove

\begin{equation}\label{eq:partial-at-34}
\boxed{F<1-(9/20)^2=319/400.}
\end{equation}

This completes all parameter and clipping cases in \(D_C\).

\paragraph{Exact optimization in width and angle}

Completing the square in \(C\) in \eqref{eq:partial-at-30} gives the unrestricted bound

\begin{equation}\label{eq:partial-at-35}
F_*\le B(\kappa):=
\frac12+\frac K2+\frac K2\sqrt{1+\kappa^2}
-\frac{\kappa(1-K)}2+\frac{\kappa^2K^2}{8}.
\end{equation}

The maximizing width for this quadratic is
\(C=1/2+K(2-\kappa)/4\); no assumption that a cap realizes the
critical supports or this width is required for the upper bound.
On \([3/8,4/5]\),

\begin{equation}\label{eq:partial-at-36}
\begin{aligned}
B'(\kappa)
&=\frac{K\kappa}{2\sqrt{1+\kappa^2}}
-\frac{1-K}{2}+\frac{\kappa K^2}{4}\\
&\le\frac{2K}{5}-\frac{1-K}{2}+\frac{K^2}{5}
=\frac{5K-3}{10}<0.
\end{aligned}
\end{equation}

Hence the largest value occurs at \(\kappa=3/8\), where

\begin{equation}\label{eq:partial-at-37}
B(3/8)=\frac{169}{512}
+\frac{K(167+16\sqrt{73})}{256}.
\end{equation}

The exact rational bounds \(K<29/70\) and
\(\sqrt{73}<171/20\) now give

\begin{equation}\label{eq:partial-at-38}
B(3/8)<\frac{169}{512}
+\frac{29(167+16\cdot171/20)}{70\cdot256}
=\frac{5259}{6400}
=\frac{411}{500}-\frac9{32000}.
\end{equation}

All other cases were bounded strictly below \(41/50\), \(4/5\),
or \(319/400\), each smaller than \(5259/6400\). The earlier
outgoing-only theorem covers the initial interval
\([\pi/4,\arctan(5/4)]\) by a value below \(41/50\).
Fixed-angle attainment and the canonical reduction therefore prove the asserted universal inequality.

For the exact reference comparison, the inequality $\arctan x\ge x-x^3/3$ for $x\ge0$, with
\(a=149/500\), gives

\begin{equation}\label{eq:partial-at-39}
\begin{gathered}
4a^3+3a-1=-\frac{4551}{31250000}<0,\\
\frac\Mc2>\frac{1+4a^2+a-a^3/3}{2}
=\frac{616648051}{750000000}>
\frac{411}{500}.
\end{gathered}
\end{equation}



\end{proof}

\subsection{A width cut with the actual partial barrier}

\begin{proposition}\label{app:partial-width}
A canonical proper joint maximum has $C<37/50$.
\end{proposition}
\begin{proof}
By \cref{lem:preliminary-cuts}, it suffices to exclude
$37/50\le C<4/5$ and $0\le h=h_L+h_R<17/50$.
Relabel the high middle endpoint as the right one, reflecting both
the cap and its actual barrier if necessary. Put $a=2C$ and let
$q_\pm$ be the values of this relabeled barrier at $\pm a/2$.
The endpoint equations and roof concavity give
\[
 e_R=\tfrac12+\tfrac h4+(3q_+-q_-)/4\le1,\qquad
 e_L=\tfrac12-\tfrac{3h}4+(3q_--q_+)/4\le1-3h/2.
\]
Use the upper bound on $e_L$ in the first equation, and the upper
bound on $e_R$ in the second, to obtain
\[
 e_R\ge\tfrac13+\tfrac h2+\tfrac23q_+,\qquad
 e_L\ge\tfrac13-\tfrac{2h}3+\tfrac23q_-.
\]
Lift the low wing by $h$, and use the diagonal parameters $u,v$
and the functions $F_h,S,G_h$ of \eqref{eq:three-tent-objective}.
Write $r=\sqrt2$, $k=r-1$, $c_8=\cos(\pi/8)$,
$\gamma=1-r/2$ and $t=2c_8-r$. The outer deficits
$w_R=a-u$, $w_L=a-v+h$ consequently satisfy
\[
 0\le w_R\le\tfrac23-\tfrac h2,\quad
 0\le w_L\le\tfrac23+\tfrac{2h}3,\quad
 w_R+w_L\le1+\tfrac h2.
\]
The last inequality uses $e_R+e_L\ge1-h/2$.

Here the three-angle calculation must retain the floor. The exterior
support triangles and the two extra tent triangles are exactly those
in the proof of \cref{lem:full-three-tent-score}: all their angles
are visited, and the actual barrier dominates their tents. A triangle
whose signed peak is $Z$ loses, when clipped by the zero floor, at
most $Z_-^2/r$. Completing the same two squares in the exterior
support deficits therefore gives
\[
 P_\alpha(U)\le G_h(a,u,v)+
 \gamma\bigl[(w_R-t)_+^2+(w_L-t)_+^2\bigr].
\]
This derivation uses no upper bound on $q_\pm$ and applies to the
relabeled actual barrier. The convex function in brackets attains
its maximum over the displayed polygon on the sum boundary.
Its endpoints are
$(2/3-h/2,1/3+h)$ and $(1/3-h/6,2/3+2h/3)$.
The second pair is more spread out and has the larger value;
since $1/3-h/6<t$, that value is $(D+2h/3)^2$, where
$D=2/3-t$.

Set $a_0=37/25$, $A=1-k$ and
$\sigma_0=(3r-1)a_0/2-(4c_8-2)$.
The supporting-plane estimate of \cref{lem:full-three-tent-critical}
now gives
\[
 \begin{gathered}
 P_\alpha(U)\le B(a_0)-kh/2-h^2/8
       -\tfrac29(\sigma_0-Ah/4)^2+\gamma(D+2h/3)^2,
 \\
 B(a)=a-(a-k)^2/2.
 \end{gathered}
\]
The derivative of the right side with respect to $h$ is
\[
 -k/2+A\sigma_0/9+4\gamma D/3
  +h(-1/4-A^2/36+8\gamma/9).
\]
The rational bounds $k>2/5$, $A<3/5$, $\gamma<3/10$,
$703/1000<\sigma_0<71/100$ and $D<47/200$ show that the
constant term is less than $-22/375$ and the coefficient of $h$
is less than $1/60$. Thus this derivative is less than
$-53/1000$ on $[0,17/50]$. Its maximum is at $h=0$.
As $k<29/70$ and $\gamma D^2<17/1000$,
\[
 P_\alpha(U)<\frac{62}{25}\frac{29}{70}-\frac{72}{625}
 -\frac29\left(\frac{703}{1000}\right)^2+\frac{17}{1000}
 =\frac{25811237}{31500000}<\frac{41}{50}<\frac\Mc2.
\]
This contradicts the reference competitor.
\end{proof}

\subsection{A strict circle margin}

\begin{proposition}\label{app:partial-circle}
For the residual data of \eqref{eq:partial-tail-data}, and
$0\le z\le w$,
\[
 a_\alpha(-C+z)-
 \left[\sqrt{1-(C-T_L)^2}-\sqrt{1-(C+z)^2}\right]
 >39c/4400.
\]
\end{proposition}
\begin{proof}
Write $g(x)=\sqrt{1-x^2}$ and $K(x)=x/g(x)$.
The wall decreases with $z$ and the circle difference increases,
so it suffices to take $z=w=c-d$. Since $T_R\ge0$ and
$T_L\le d$, the gap there is at least
\[
 \frac cs(2C-c)-[g(C-d)-g(C+c-d)].
\]
Also $d\ge\delta=c/(1+s)$, since $H\le1$.
The bracket decreases with $d$. Put
\[
 \eta=\delta-c/2=\frac{c^3}{2(1+s)^2}\ge c^3/8,
 \qquad \mathcal A_q(y)=\frac1q
       \int_{y-q/2}^{y+q/2}K(x)\,\dd x.
\]
The gap divided by $c$ is at least
$(2C-c)/s-\mathcal A_c(C-\eta)$.
On all these intervals $K$ is increasing and convex. Hence the
average increases with its center and width and is convex in its
center. For width monotonicity, pair the positive and negative
integration variables after rescaling to $[-1/2,1/2]$; their
derivative has the sign of $K'$ increasing. A linear function of
$C$ minus a fixed-width average is concave, so its minimum on
$[1/2,37/50]$ occurs at an endpoint.

If $c\le1/3$, then
$\mathcal A_c(C-\eta)\le\mathcal A_{1/3}(C)$.
The derivative of $(2C-c)/\sqrt{1-c^2}$ with respect to $c$ is
$(2Cc-1)/(1-c^2)^{3/2}<0$. Using $\sqrt2<99/70$ gives
\[
 (2C-c)/s>\tfrac{35}{33}(2C-1/3).
\]
At $C=1/2$, the average is $2\sqrt2-\sqrt5<23/35<2/3$;
at $C=37/50$ it is $(8\sqrt{59}-\sqrt{1001})/25<6/5$.
The respective gaps exceed $4/99$ and $8/495$.
The radical comparisons follow by squaring
$\sqrt2<10/7$, $\sqrt5>11/5$, $\sqrt{59}<77/10$ and
$\sqrt{1001}>158/5$. Concavity proves a gap greater than
$8/495>39/4400$ throughout.

If $1/3\le c<3/\sqrt{73}$, take $q=44/125$.
Then $\eta\ge1/216$, $c<q$ and
$s\le2\sqrt2/3<33/35$, so
\[
 \mathcal A_c(C-\eta)\le\mathcal A_q(C-1/216),\qquad
 (2C-c)/s>\tfrac{35}{33}(2C-44/125).
\]
At $C=1/2$ the integration endpoints are $8623/27000$ and
$18127/27000$. Their $g$ values are respectively below $19/20$
and above $37/50$, so the average is below $105/176<2/3$.
The resulting gap exceeds $17/825$.
At $C=37/50$ the endpoints are $15103/27000$ and $24607/27000$.
Their $g$ values are below $829/1000$ and above $411/1000$:
the exact squared comparisons are
\[
 15103^2=228100609>228001311
       =27000^2[1-(829/1000)^2],
\]
\[
 24607^2=605504449<605856591
       =27000^2[1-(411/1000)^2].
\]
Thus the average is below $19/16$, and its gap exceeds
$329/275-19/16=39/4400$. Concavity proves this strict lower
bound throughout the second interval, proving the proposition.
\end{proof}

\subsection{Reflected supports and the terminal impulse}

\begin{proposition}\label{app:partial-reflection}
For a residual canonical proper joint maximum, the error in
\eqref{eq:partial-weighted-error} satisfies
\[
 \mathcal E_\theta=0\quad(0<\cot\alpha\le1/8),\qquad
 \mathcal E_\theta<6517/6400000
       \quad(1/8\le\cot\alpha<3/8).
\]
\end{proposition}
\begin{proof}
For a rising middle use the left-wing companion surrogate of
\cref{lem:partial-moment}. For a descending middle use the
right-wing first-support surrogate: it equals the actual first
support on every visited angle and removes the unused middle atom.
Saturation gives its unused tail at $(C,1-h_R)$.
In every case the companion unused tail is
$G(t)=2C\sin t+e_L\cos t$ for $\alpha<t<L$.
Set $\varepsilon=L-\alpha$, $\delta=c/(1+s)$ and reflect the
support equations, without reflecting their spatial objective:
\[
 P(r)=-q(L-r),\quad Q(r)=-p(L-r),\quad
 U(r)=v(L-r),\quad V(r)=u(L-r).
\]
The curvature and shadow laws give
\begin{equation}\label{eq:partial-reflected-controls}
 \begin{gathered}
 P'=U-1-Q,\quad Q'=V-1+P,\quad P\le1,\ Q\ge-1,\\
 0\le U\le\kappa(Q),\quad0\le V\le\kappa(P),\\
 P,Q>0\Rightarrow U=0,\ V\le1/2,\qquad
 P,Q<0\Rightarrow V=0,\ U\le1/2.
 \end{gathered}
\end{equation}
In the final removed interval of the rising surrogate both
curvatures are zero directly; no arm law is needed there.
The initial values are $(P(0),Q(0))=(e,D-1)$, where
\[
 \begin{aligned}
 (e,D)=(e_L,3C+T_R)\quad\hbox{for rising or horizontal middle},
 \\
 (e,D)=(e_L+h_R,3C)\quad\hbox{for descending middle}.
 \end{aligned}
\]
Both densities vanish on $(0,\varepsilon)$.
At $\varepsilon$, $P$ is continuous and $Q$ jumps upwards by
$j=\tau/s\le\delta$. There is no companion terminal atom or
later visited atom. The free initial trajectory is
\[
 P=1+(e-1)\cos r-D\sin r,\qquad
 Q+1=D\cos r+(e-1)\sin r.
\]

We first record a common energy and duration calculation.
Put $\mathcal H=(P-1/2)^2+(Q+1)^2$.
Expansion at the impulse gives
\[
 \mathcal H(\varepsilon+)-\mathcal H(0)
 =(1-e)(1-s)-Dc+2j[Ds-(1-e)c]+j^2.
\]
In the parameter boxes below $Ds-(1-e)c>1$; this expression
increases with $j\ge0$. At $j=\delta$ it equals
\begin{equation}\label{eq:partial-impulse-energy}
 \delta^2\left[\frac{2e-1+\delta^2}{1+\delta^2}-Dc\right].
\end{equation}
After the impulse, in the initial $P,Q>0$ sector,
$\mathcal H'=2(Q+1)(V-1/2)\le0$.
If the impulse cost is at most $\eta_0$, then at the first
$P$-zero, or at $\varepsilon$ if that zero occurred earlier,
\[
 Q+1\le\sqrt Z,\qquad Z=D^2-e(1-e)+\eta_0.
\]
The initial $Q$ component remains positive across the unused gap.
Every later positive $Q$ component has amplitude at most $1/8$:
it starts at $Q=0,P\le1$; while $P>0$, the shadow laws give
$P'\le-1$, $Q'\le P-1/2$, so its total positive increase is
at most $\int_0^{1/2}(1/2-r)\,\dd r=1/8$.
Once $P\le1/2$, $Q'$ is nonpositive, and once $P\le0$ it is
strictly negative. $P$ cannot rise while $Q>0$.
Thus only the initial component can have $U>1$, since that
requires $Q>1$ by \eqref{eq:partial-reflected-controls}.

After $P\le0$ and while $Q>1$, write $\rho$ for elapsed time.
Then $P'\le-1$, hence $P\le-\rho$ and, for $\rho\le1$,
$Q'\le-1/2-\rho/2$. For $-1\le P\le0$ use
$V\le(1-P)/2$; for $P<-1$ use $V\le-P$, giving $Q'\le-1$.
Consequently
\begin{equation}\label{eq:partial-excess-decay}
 Q-1\le\sqrt Z-2-\rho/2-\rho^2/4,
 \qquad D_0(Z)=\sqrt{1+4(\sqrt Z-2)_+}-1.
\end{equation}
All the ensuing bounds on $D_0$ are less than one, closing this
comparison. Once $0<Q\le1$, one has $P'\le-(1+Q)/2$ and
$Q'\le-1/2$, so this component cannot return to $Q>1$.
While $P,Q>0$, $Q'\ge-1$ and the impulse is upward;
therefore $P(r)\le e-Dr+r^2/2$. Its first zero, if the
$Q$ component has not already ended, occurs no later than
\[
 \tau_0(D,e)=D-\sqrt{D^2-2e}.
\]
The last possible excess time is thus at most
$\max(\varepsilon,\tau_0)+D_0(Z)$.

\paragraph{The near-full range.}
If $\cot\alpha\le1/8$, then
$\varepsilon<1/8$, $\delta<1/16$, and $s>124/125$.
The ordinary visited niche has endpoint box
$H_0(C)=1-\sqrt{1-C^2}<1/3$.
The outgoing endpoint is nonpositive at $C$ by top projection,
and at $-C$ is at most $2C\cot\alpha<H_0(C)$:
$H_0(C)\ge C^2/2>C/4$. Thus the actual barrier endpoints
are also below $1/3$. Endpoint pressure gives
\[
 \begin{array}{c|cc}
 &e&D\\\hline
 \text{rising/horizontal}&0\le e<3/4&3/2<D<1801/800\\
 \text{descending}&0\le e<19/25&3/2<D<111/50.
 \end{array}
\]
For the second row use $h_R<1-s<1/125$; for the first use
$T_R\le\delta/2<1/32$.
Also $Ds-(1-e)c>(3/2)(124/125)-1/8>1$.
The bracket in \eqref{eq:partial-impulse-energy} is below
$53/100$, so its cost is below $1/480$.
Take $\eta_0=1/480$.

If $e\le3/10$ and $Z>4$, then $D>1999/1000$, and
$\tau_0<4/25$, since
$2(1999/1000)(4/25)-(4/25)^2-3/5=44/3125>0$.
Uniformly $\sqrt Z<451/200$ and $D_0<17/40$.
Thus the last excess precedes $117/200<18/25$.
If $Z\le4$ there is no excess after the zero or impulse.

If $e\ge3/10$, then $P(\varepsilon)\ge e-Dc>0$.
The zero follows the impulse and the last excess is bounded by
$W=\tau_0+D_0(Z)$. On $Z>4$ put
$A=\sqrt{D^2-2e}$, $S=\sqrt Z$, $Q_0=\sqrt{4S-7}$.
Then
\begin{equation}\label{eq:partial-duration-derivatives}
 W_e=1/A+(2e-1)/(Q_0S),\qquad
 W_D=1-D/A+2D/(Q_0S).
\end{equation}
Here $D>2$, $D/A<4/3$, $S<9/4$, $Q_0<3/2$.
When $e<1/2$, the negative term in $W_e$ is at most $1/5$,
whereas $1/A>11/25$; when $e\ge1/2$ positivity is immediate.
Also $W_D>-1/3+32/27>0$.
The expression clipped by $D_0=0$ remains increasing in $e$
across $Z=4$. If increasing $e$ to its upper endpoint gives
$Z\le4$, its value is below $9/20$. Otherwise increase also
$D$ and check the two endpoint pairs:
\[
 \begin{array}{c|ccc}
 (D,e)&\tau_0&\sqrt Z&D_0\\\hline
 (1801/800,3/4)&<109/300&<221/100&<107/300\\
 (111/50,19/25)&<3/8&<109/50&<5/16.
 \end{array}
\]
All entries follow by squaring positive rationals. For the first
row the squared margins are $1403/360000$, $2669/1920000$
and $49/90000$; the second row gives $259/60000$ for the
middle entry and $4(109/50)-7=43/25<(21/16)^2$ for the last.
The sums are below $18/25$ and $11/16$ respectively.
It follows that $U(r)\le1$ for almost every $r\ge18/25$.
The Taylor bound $\cos x\ge1-x^2/2+x^4/25$ on $[0,1]$
gives $\cos(18/25)>37/50$. Therefore
$v(t)\le1$ on $[0,\arcsin C]$, which contains $[0,\theta]$.
This proves $\mathcal E_\theta=0$ in the near-full range.

\paragraph{The complementary range.}
Now $1/8\le\cot\alpha<3/8$. Then $c<3/8$, $s>14/15$,
$\delta<2/11$ and $\varepsilon<9/25$.
The last inequality follows from
$\arctan(3/8)<3/8-(3/8)^3/3+(3/8)^5/5<9/25$.
For rising or horizontal middle, top projection bounds the actual
left outgoing endpoint by
\[
 \cot\alpha(2C-\delta/2)
 =2C\cot\alpha-(\sqrt{1+\cot^2\alpha}-1)/2<419/800.
\]
This expression increases with both $C$ and $\cot\alpha$;
evaluate at $37/50,3/8$ and use $\sqrt{73}>17/2$.
For descending middle it is below $2C\cot\alpha<111/200$,
and $h_R<1/15$. With the ordinary endpoint box below $1/3$,
the actual pressure equations yield
\[
 \begin{array}{c|cc}
 &e&D\\\hline
 \text{rising/horizontal}&0\le e<179/200&3/2<D<2311/1000\\
 \text{descending}&0\le e<1&3/2<D<111/50.
 \end{array}
\]
Indeed the first pressure is below $2857/3200$, and
$3C+T_R<111/50+1/11=1271/550<2311/1000$.
The second pressure is below $1/2+1/12+333/800=2399/2400$.
Now $Ds-(1-e)c>(3/2)(14/15)-3/8>1$.
The energy cost is below $\delta^2<1/30$, so take
$\eta_0=1/30$.
Uniformly $\sqrt Z<58/25$ and $D_0(Z)<13/25$.
If the first zero precedes the impulse, the episode ends before
$9/25+13/25=22/25<91/100$.
Otherwise use $W=\tau_0+D_0$ and
\eqref{eq:partial-duration-derivatives} on $Z>4$.
Both derivatives are positive for $0\le e\le1$:
$S^2-A^2=e+e^2+1/30>0$ and $Q_0\ge1$ control the possible
negative term in $W_e$; also $D^2>119/30$,
$D/A<\sqrt{119/59}<3/2$, $S<7/3$, $Q_0<31/20$, so
$W_D>-1/2+180/217>0$.
Across $Z=4$ the clipped expression remains increasing in $e$.
If its upper-$e$ value has $Z\le4$, then $\tau_0(D,1)<3/5$
because $D^2>119/30>(59/30)^2$. Otherwise increase $D$ too.
The two endpoint certificates are
\[
 \begin{array}{c|ccc}
 (D,e)&\tau_0&\sqrt Z&D_0\\\hline
 (2311/1000,179/200)&<107/250&<2299/1000&<241/500\\
 (111/50,1)&<51/100&<223/100&<39/100.
 \end{array}
\]
For the first row the three squared margins are $629/125000$,
$3193/600000$ and $81/250000$; for the second they are
$43/10000$, $67/6000$ and $121/10000$.
The sums are below $91/100$ and $9/10$ respectively.

On an excess episode the polynomial in
\eqref{eq:partial-excess-decay} factors as
$(D_0-\rho)(1/2+(D_0+\rho)/4)$.
For $0\le\rho\le D_0<13/25$ its second factor is below $19/25$.
Since $U-1\le Q-1$ whenever $U>1$, we conclude
\[
 (U(r)-1)_+\le\tfrac{19}{25}(91/100-r)_+,
 \quad (v(t)-1)_+\le\tfrac{19}{25}(t-t_0)_+,
 \quad t_0=L-91/100.
\]
Finally $\cos\theta>2/3$ and
$\theta-t_0<7/40$. The latter follows from
$\cos(147/200)>37/50$ using the same Taylor lower bound;
its rational margin is $62448881/40000000000$.
If $\theta\le t_0$ the error is zero. Otherwise
\[
 \mathcal E_\theta\le
 \frac{19}{25\cos\theta}\int_{t_0}^{\theta}
          (t-t_0)(\theta-t)\,\dd t
 <\frac{19}{100}(7/40)^3=\frac{6517}{6400000},
\]
because its exact kernel is
$\lambda\cos t-\sin t=\sin(\theta-t)/\cos\theta$.
As $c\ge1/\sqrt{65}>3/25$, this is less than
$117/110000<39c/4400$, as required.
\end{proof}
